% Generated from the matching Typst scaffold by
% scripts/migrate_typst_report_to_latex.py (prose drafted 2026-06-25).

% ── 2.3  Cognitive Foundations of Forgetting ─────────────────────────────────
% PURPOSE: Ground the two principled policies (Type-Aware Decay, CLS Consolidation) in
%          cognitive science; transience as adaptive rather than as failure.
% ─────────────────────────────────────────────────────────────────────────────

\section{Cognitive Foundations of Forgetting}

Rather than treat forgetting as loss, cognitive science treats it as adaptive --- a memory
system that discards on a rational schedule generalises and decides better than one that
retains indiscriminately~\cite{richards2017}. Two principles from that literature ground
the two content-aware policies this thesis evaluates.

The first is the Anderson--Schooler rational analysis of memory~\cite{anderson1991}. They
show that the probability a memory will be \emph{needed} again falls predictably with the
time since it was last used and rises with how often it has been retrieved --- so an
optimal memory system should let need-probability, not raw age, govern retention. This is
the direct basis for \textbf{Type-Aware Decay}, whose retention score is the multiplicative
form $\text{base} \times \text{age}^{-d(\text{type})} \times (1 + \text{retrieval\_count})^{0.5}$:
age decays the score, retrieval reinforces it, and the per-type exponent $d(\text{type})$
lets different kinds of record age at different rates
(Section~\ref{sec:type-decay}).

The second is the Complementary Learning Systems theory~\cite{mcclelland1995}: the brain
pairs a fast-learning hippocampal store with a slow-learning neocortical one, and
\emph{consolidation} gradually transfers and summarises the stable structure from the
former into the latter. This grounds \textbf{CLS Consolidation}, which on a fixed schedule
(every $k=5$ tasks) clusters older memories and replaces each cluster with an
LLM-generated summary, compressing repeated structure rather than storing every instance
verbatim (Section~\ref{sec:cls}).

Both policies treat transience as a feature. Crucially, the parameters that operationalise
the cognitive theory --- in particular the per-type decay exponents $d(\text{type})$ ---
are frozen at design time from the theoretical values, \emph{not} calibrated from pilot
data, so that the policies test the cognitive hypothesis rather than a curve fitted to the
benchmark. The next two sections of this background motivate why retrieval, not just
storage, shapes whether any of this helps; the operational definitions follow in
Chapter~\ref{ch:method}.
